\part{Práctica 1: Contexto y Requisitos}

\section{Introducción}
El sector de la restauración ha experimentado una digitalización acelerada, mejorando la interacción entre usuarios y establecimientos mediante plataformas de reserva y descubrimiento gastronómico \cite{Gomez2023}. Sin embargo, la gestión del aforo y la disponibilidad en tiempo real sigue siendo un reto operativo. Este proyecto propone un sistema integral que permite a los restaurantes registrarse, generar códigos QR únicos y controlar su aforo en tiempo real mediante el escaneo por parte de los clientes, facilitando así un registro preciso de entradas, salidas y acompañantes \cite{TechRestaurant2022}.

\section{Definición del Problema}
El problema central radica en la ``Crisis de Información y Ambigüedad'' que sufren los clientes al intentar acudir a un restaurante sin reserva previa durante picos de afluencia. La dependencia de la intervención humana (el hostelero actualizando su estado manualmente en una app) genera datos desfasados e inexactos. 

La solución que se pretende implementar es la automatización del control de aforo mediante tecnología QR, eliminando el error humano y garantizando que el dato de ocupación sea exacto y transaccional en tiempo real.

\section{Alcance del Sistema}
El sistema abarcará el desarrollo de una plataforma web dual (interfaz para restaurantes e interfaz para clientes). Incluirá la generación y escaneo de códigos QR, control transaccional de aforo (entradas y salidas), cálculo de métricas (historial de afluencia y duración media) y un sistema de valoraciones. 

Fuera del alcance quedarán, en esta fase inicial, pasarelas de pago, gestión de comandas internas del restaurante o integración con hardware físico de barreras/tornos.

\section{Objetivos del Sistema}
\begin{itemize}
    \item Desarrollar un sistema web de control de acceso basado en códigos QR para restaurantes.
    \item Implementar un motor de generación y validación de códigos QR únicos por establecimiento.
    \item Automatizar el control de aforo en tiempo real mediante el escaneo de dichos códigos (registrando cliente y acompañantes).
    \item Proporcionar al usuario final un listado de restaurantes filtrable (por cercanía y tipo de comida) con disponibilidad 100\% real.
    \item Dotar al restaurante de métricas e historial de afluencia y duración media de las visitas.
\end{itemize}

\section{Antecedentes}
Actualmente, plataformas consolidadas como TheFork dominan el mercado. Un análisis de usabilidad de dichas plataformas revela aspectos muy positivos, como una interfaz intuitiva, una excelente gestión de reservas a futuro y una alta utilidad percibida por el usuario \cite{TheFork2023}. No obstante, presentan deficiencias críticas en la operativa diaria:

\begin{itemize}
    \item \textbf{Falta de visibilidad en ``Horas Punta'':} El sistema falla al no indicar claramente si el aforo está completo en el momento exacto de la consulta.
    \item \textbf{Información desfasada:} La disponibilidad mostrada depende de la actualización manual constante del hostelero, lo cual es un error frecuente que penaliza la usabilidad.
    \item \textbf{Ambigüedad:} Al no mostrar el estado del local en tiempo real, se genera incertidumbre en el usuario.
\end{itemize}

\section{Extracción de Requisitos}

\subsection{Captura de Requisitos (Elicitation)}
Para la captura de requisitos se ha realizado un análisis de sistemas existentes, específicamente la plataforma \textit{TheFork}. De esta entrevista indirecta con las necesidades del mercado se han extraído los siguientes indicios:
\begin{itemize}
    \item Necesidad de visibilidad en tiempo real para evitar la incertidumbre en horas punta.
    \item Automatización del flujo de datos para evitar información desfasada por error humano.
\end{itemize}

\subsection{Definición de Requisitos}
Siguiendo la metodología UWE, se clasifican los requisitos en tres categorías fundamentales:

\subsubsection{Requisitos de Información (RI)}
\begin{itemize}
    \item \textbf{RI-01:} Datos del restaurante (Nombre, dirección, tipo de cocina, aforo máximo).
    \item \textbf{RI-02:} Registro de afluencia (Historial de entradas/salidas con fecha y hora).
\end{itemize}

\subsubsection{Requisitos Funcionales (RF)}
A continuación se listan los requisitos funcionales en lenguaje natural:
\begin{itemize}
    \item \textbf{RF-01:} Registro y gestión de perfil para establecimientos.
    \item \textbf{RF-02:} Generación de código QR único vinculado al ID del restaurante.
    \item \textbf{RF-03:} Registro de entrada mediante escaneo de QR.
    \item \textbf{RF-04:} Control de aforo automatizado.
\end{itemize}

\subsubsection{Requisitos No Funcionales (RNF)}
\begin{itemize}
    \item \textbf{RNF-01:} Interfaz \textit{responsive} optimizada para dispositivos móviles.
    \item \textbf{RNF-02:} Seguridad y firma digital en los códigos QR.
    \item \textbf{RNF-03:} Tiempo de respuesta del servidor inferior a 3 segundos.
\end{itemize}

\subsection{Validación de Requisitos}
Se ha aplicado la técnica de \textbf{Walkthrough} (Revisión formal) \cite{Sommerville2015}. El equipo (en este caso el autor) ha simulado el recorrido de los casos de uso principales para detectar inconsistencias, confirmando que el catálogo de requisitos cubre el 100\% de los objetivos del sistema planteados en la Práctica 1.
